% pccx — arXiv-ready preprint scaffold.
%
% Goal: a single self-contained LaTeX file that builds with `latexmk
% -pdf main.tex` using any stock TeX Live distribution.  Drop in
% content as sections close over time; the file is designed so a
% weekend build can produce a v0.1 draft.
%
% Compile:   latexmk -pdf main.tex
% Output:    main.pdf  (keep transient artefacts out of git via
%                       ../../.gitignore's LaTeX block)
%
% Bibliography:
%   refs.bib co-located with main.tex.  Sphinx's refs.bib is the
%   canonical source — keep entries in sync manually until a
%   generator lands.

\documentclass[11pt,a4paper]{article}

%-- Encoding & fonts -------------------------------------------------
\usepackage[T1]{fontenc}
\usepackage[utf8]{inputenc}
\usepackage{lmodern}
\usepackage{microtype}

%-- Layout ----------------------------------------------------------
\usepackage[margin=1in]{geometry}
\usepackage{parskip}
\usepackage{titlesec}
\titleformat{\section}{\large\bfseries}{\thesection}{0.6em}{}
\titleformat{\subsection}{\normalsize\bfseries}{\thesubsection}{0.6em}{}

%-- Floats / listings / code ---------------------------------------
\usepackage{graphicx}
\usepackage{booktabs}
\usepackage{listings}
\lstset{
  basicstyle=\ttfamily\small,
  breaklines=true,
  columns=fullflexible,
  xleftmargin=1em,
  frame=none,
  aboveskip=0.6em,
  belowskip=0.6em,
}

%-- Math -----------------------------------------------------------
\usepackage{amsmath, amssymb, amsthm}

%-- References -----------------------------------------------------
\usepackage[numbers,sort&compress]{natbib}
\usepackage[dvipsnames]{xcolor}
\usepackage{hyperref}
\hypersetup{
  colorlinks=true,
  linkcolor=blue!50!black,
  citecolor=blue!50!black,
  urlcolor=blue!50!black,
}
\usepackage{url}

%-- Authorship -----------------------------------------------------
\title{%
  pccx: An Open, Sail-Specified W4A8 NPU for\\
  Edge LLM Inference on the Xilinx Kria KV260
}
\author{%
  Hyunwoo Kim\\
  \texttt{k1h6w4@gmail.com}\\
  \small \url{https://hwkim-dev.github.io/pccx/}
}
\date{Draft v0.1 --- \today}

\begin{document}
\maketitle

\begin{abstract}
We present \textbf{pccx}, an open, scalable Neural Processing Unit
(NPU) architecture for autoregressive decoding of Transformer-based
large language models on resource-constrained edge devices.  The v002
reference implementation targets the Xilinx Kria KV260 (Zynq
UltraScale+ ZU5EV), is formally specified in the
Sail~\cite{armstrong2019isa} ISA semantics language (the same language
used for RISC-V, Arm, CHERI and Morello), and implements recent W4A8KV4
quantisation~\cite{qserve2024}, FlashAttention~\cite{dao2023flash2},
and NPU-adapted speculative decoding~\cite{openpangu2026} in a single
64-bit VLIW ISA with five base opcodes (GEMV / GEMM / MEMCPY /
MEMSET / CVO).  An accompanying desktop profiler, \textbf{pccx-lab},
supplies 16 built-in trace analyzers and an AI-driven UVM strategy
copilot.  This preprint covers the architecture, the Sail model, the
profiler design, and the current measurement status.
\end{abstract}

\section{Introduction}
\label{sec:introduction}

\textbf{TODO} — three paragraphs:
\begin{itemize}
  \item Edge LLM inference landscape (Gemma 3N, Llama, Phi).
  \item Why W4A8KV4 on FPGA makes sense for a KV260-class device.
  \item pccx contribution summary; pointer to the three repos at
        \url{https://github.com/hwkim-dev/}.
\end{itemize}

\section{Architecture}
\label{sec:architecture}

\subsection{Compute cores}

32$\times$32 systolic GEMM array with cascade break at row 16;
4-lane $\times$ 32-MAC GEMV with 5-stage reduction; single BF16 SFU
pipeline covering EXP / SQRT / GELU / SIN / COS / REDUCE\_SUM /
SCALE / RECIP.  Weight stationary in GEMM, weight streaming in GEMV.

\subsection{Memory hierarchy}

1.75~MB L2 URAM at the geometric centre of the floorplan.  Dual-port
access: ACP for latency-critical host path, NPU port-B for direct
compute fills.  Constant cache with 64 entries, 6-bit pointer from
the ISA.  HP0/1 carry GEMM weights, HP2/3 carry GEMV weights; both
pairs independent at 128-bit $\times$ 250 MHz.

\subsection{ISA}

64-bit VLIW with a 4-bit opcode in bits \texttt{[63:60]} and a 60-bit
body.  Five opcodes:

\begin{center}
\begin{tabular}{llp{0.5\textwidth}}
\toprule
Opcode   & Encoding & Body layout \\
\midrule
\texttt{GEMV}   & \texttt{0x0} & dst(17) $\|$ src(17) $\|$ flags(6) $\|$ size(6) $\|$ shape(6) $\|$ lanes(5) $\|$ rsv(3) \\
\texttt{GEMM}   & \texttt{0x1} & identical to \texttt{GEMV}  \\
\texttt{MEMCPY} & \texttt{0x2} & from(1) $\|$ to(1) $\|$ dst(17) $\|$ src(17) $\|$ aux(17) $\|$ shape(6) $\|$ async(1) \\
\texttt{MEMSET} & \texttt{0x3} & cache(2) $\|$ dst(6) $\|$ a(16) $\|$ b(16) $\|$ c(16) $\|$ rsv(4) \\
\texttt{CVO}    & \texttt{0x4} & func(4) $\|$ src(17) $\|$ dst(17) $\|$ len(16) $\|$ flags(5) $\|$ async(1) \\
\bottomrule
\end{tabular}
\end{center}

\section{Sail Formal Specification}
\label{sec:sail}

The ISA is authored in \textbf{Sail}~\cite{armstrong2019isa}, with
modules split incrementally per the language conventions:

\begin{lstlisting}[language={}]
formal/sail/
  pccx.sail_project
  src/prelude.sail       # base bitvectors
  src/pccx_types.sail    # opcode table, body structs, flags
  src/pccx_regs.sail     # cycle / pc / committed_any
  src/pccx_decode.sail   # 64-bit word -> typed instr union
  src/pccx_execute.sail  # step(); advance_cycle; record_event()
  tests/smoke_decode.sail
\end{lstlisting}

Every SystemVerilog \texttt{typedef} in the RTL ISA package has a
1:1 Sail counterpart, so width drift fails the Sail type checker
before the synthesiser.

\section{pccx-lab Profiler}
\label{sec:lab}

\textbf{TODO} — describe:
\begin{itemize}
  \item 16 built-in \texttt{TraceAnalyzer} implementations.
  \item \texttt{Copilot::investigate} / \texttt{suggest\_fix} surface.
  \item Research-citation registry (21 papers auto-cited).
  \item Tauri desktop UI with WaveformViewer, FlameGraph, ISA replay.
\end{itemize}

\section{Methodology}
\label{sec:method}

Two sample traces are bundled with the repository
(\texttt{samples/}, 16- and 128-token decode) so every number in this
paper is reproducible.  See the Quickstart at
\url{https://hwkim-dev.github.io/pccx/en/docs/quickstart.html}.

\section{Evidence}
\label{sec:evidence}

\textbf{TODO} — fill as measurements land:
\begin{itemize}
  \item Gemma-3N E4B tokens per second on KV260.
  \item KV260 resource usage (LUT / DSP / URAM / BRAM36).
  \item Timing closure @ 400~MHz core / 250~MHz AXI.
  \item CPU baseline (Ryzen 4500U, llama.cpp Q4\_K\_M).
  \item GPU baseline (RTX 4060, HF Transformers BF16).
\end{itemize}

Current status is tracked live at
\url{https://hwkim-dev.github.io/pccx/en/docs/Evidence/index.html}.

\section{Related Work}
\label{sec:related}

Sail precedents for ISA formalisation: \cite{armstrong2019isa} for
RISC-V; \cite{sailarm2020} for Armv8-A.  W4A8KV4 lineage:
\cite{qserve2024} (QServe), \cite{qqq2024} (QQQ), \cite{sramfreq2025}
(SRAM--frequency bottleneck model).  FlashAttention:
\cite{dao2023flash2} (FA-2), \cite{dao2024flash3} (FA-3).
Speculative decoding on static-graph NPUs: \cite{openpangu2026}.

\section{Conclusion and Future Work}
\label{sec:conclusion}

\textbf{TODO} — one paragraph: the value of combining open RTL +
formal Sail spec + profiler; near-term goal of board-captured
tokens per second; longer-term goals (Gemma-4, Isabelle proofs).

\bibliographystyle{plainnat}
\bibliography{refs}

\end{document}
